\section{Conclusion}

Parmi les approches étudiées,
l'utilisation d'une file de priorité (PQ) combinée à une recherche par préfixe est la plus efficace pour traiter les requêtes.

Les méthodes basées sur le tri sont moins performantes,
car elles nécessitent de trier l'ensemble des mots correspondants,
ce qui entraîne un coût plus élevé.

L'approche avec une file de priorité permet de sélectionner les meilleurs
éléments sans effectuer un tri complet. Elle traite les résultats en $O(m \log \min(k, m))$,
ce qui évite de trier les $m$ éléments.
Cette méthode est particulièrement adaptée au contexte où le nombre de résultats à afficher $k$
est généralement bien inférieur au nombre de termes correspondants $m$.
